\section{Step 64 Quantize and retain the ordinary count}\label{h263-block-s04}
Build the quantizer request from the private forward block, parameter, block type
and mode. Its output supplies the selected-level and reconstructed-coefficient
arrays. Save the returned ordinary nonzero count in private storage before
calling the inverse, which uses the same return registers for its own result.

The count excludes INTRA DC. It is useful for later coding decisions, but it
is not a six-block coded-pattern mask and does not yet describe the emitted bits.

\paragraph{Observe} The default INTER block selects ordinary DC level one and
returns ordinary nonzero count one.
\Needspace{8\baselineskip}
\paragraph{Check} Why save the quantizer count across the next call?
\begin{itemize}
  \item[A.] The inverse call can replace the return registers.
  \item[B.] It is the picture's byte stride.
  \item[C.] The count changes the current source bytes in place.
\end{itemize}
